% SPDX-License-Identifier: GPL-3.0-or-later
% Copyright (C) 2026 Alessandro Greco
\chapter{Дополнительная криптография}\label[appendice]{cap:B-crittografia}
\aggiornatoal{23 сентября 2026 г.}{dir-spec и tor-spec (коммит torspec \texttt{8e1f245}); библиотека \texttt{cryptography}; лабораторная сеть}

\Cref{cap:05-crittografia} вводила криптографические основы Tor с минимальной математической подготовкой. В этом приложении подробнее рассматриваются три ранее лишь кратко упомянутые темы, каждая с воспроизводимыми примерами: ключи X25519 и принудительная фиксация отдельных битов, преобразование между формами Монтгомери и Эдвардса одной кривой с проверкой по реальному сертификату, а также изменяемость потокового шифртекста, лежащая в основе атаки маркировки (\enquote{marking attack}) из \cref{sec:tor1-debolezze}.

% ---------------------------------------------------------------------------
\section{Одна кривая, две формы}\label{sec:B-curve}

X25519 и Ed25519 используют \emph{одну и ту же} эллиптическую кривую над полем целых по модулю $p = 2^{255} - 19$, записанную в двух формах:
\begin{description}
  \item[Форма Монтгомери] $v^2 = u^3 + 486662\,u^2 + u$. В X25519 открытый ключ содержит лишь координату $u$ точки, записанную в 32 байтах с младшим байтом первым (little-endian). Эта форма допускает простой и быстрый алгоритм умножения точки на скаляр — \enquote{лестницу Монтгомери}, время выполнения которого не зависит от секретного ключа.
  \item[Скрученная форма Эдвардса] $-x^2 + y^2 = 1 + d\,x^2y^2$, где $d = -121665/121666$. В Ed25519 эта форма используется для цифровых подписей: открытый ключ кодирует координату $y$, а знак координаты $x$ хранится в старшем бите последнего байта.
\end{description}
Две формы связаны явным преобразованием; оно используется в \ref{sec:B-conversione}.

\subsection{Фиксация битов ключей X25519}

Закрытый ключ X25519 задаётся 32 байтами, но перед использованием реализация фиксирует отдельные биты: обнуляет три младших бита, обнуляет старший бит и устанавливает предпоследний в единицу. Благодаря обнулению трёх младших битов скаляр становится кратным 8 — \emph{кофактору} кривой, — поэтому компоненты малого порядка принятой точки не раскрывают сведения о скаляре. Фиксация старших битов также задаёт постоянное число шагов лестницы Монтгомери. Это можно проверить библиотекой, используемой в лабораторных работах:
\begin{lstlisting}[language=Python,numbers=none]
from cryptography.hazmat.primitives.asymmetric.x25519 import X25519PrivateKey
from cryptography.hazmat.primitives import serialization as s
import os
def public_key(k):
    return X25519PrivateKey.from_private_bytes(k).public_key().public_bytes(
        s.Encoding.Raw, s.PublicFormat.Raw)
a = bytearray(os.urandom(32)); b = bytearray(a)
b[0] ^= 0x07                     # меняются только три младших бита
print(public_key(bytes(a)) == public_key(bytes(b)))      # True
\end{lstlisting}
Две последовательности байтов закрытого ключа, различающиеся только битами, которые перезаписываются при фиксации, дают один и тот же открытый ключ.

\subsection{Точка на бесконечности}

Во время рукопожатия ntor и клиент, и ретранслятор проверяют, что ни одна операция $\mathrm{EXP}()$ не возвращает точку на бесконечности. Спецификация отмечает, что для Curve25519 эта проверка достаточно заменяет валидацию принадлежности $Y$ группе~\spec{tor-spec}{The "ntor" handshake}. Без неё атакующий мог бы передать открытый ключ малого порядка и получить предсказуемый общий секрет, например строку из нулей. Библиотека \texttt{cryptography} выполняет эту проверку самостоятельно: вызов \texttt{exchange} с полностью нулевым 32-байтовым открытым ключом возбуждает исключение (\texttt{Error computing shared key}), а не возвращает нулевой секрет.

% ---------------------------------------------------------------------------
\section{Из формы Монтгомери в форму Эдвардса: перекрёстный сертификат ntor}\label{sec:B-conversione}

Каждый ретранслятор публикует в дескрипторе открытый ключ ntor (X25519) и доказывает владение соответствующим закрытым ключом с помощью поля \texttt{ntor-onion-key-crosscert}. Оно содержит сертификат Ed25519 типа \texttt{0a}, удостоверяющий основную идентичность Ed25519 ретранслятора, но \emph{подписанный ключом ntor}~\spec{dir-spec}{ntor-onion-key-crosscert}. Для этой подписи ключ ntor необходимо преобразовать из формы Монтгомери в форму Эдвардса. Спецификация задаёт преобразование~\spec{dir-spec}{Converting a curve25519 public key to an ed25519 public key}:
\[
y = \frac{u - 1}{u + 1} \bmod p .
\]
Координата $u$ определяет $y$, но не знак $x$: одной координате Монтгомери соответствуют два симметричных открытых ключа Эдвардса. Восстановить знак по открытому ключу Монтгомери невозможно, поэтому дескриптор передаёт его отдельно в первом аргументе поля перекрёстного сертификата: 0 означает положительный, 1 — отрицательный знак. Ретранслятор вычисляет его из закрытого ключа.

Модуль \file{labs/ru/код/torlab/conversion.py} реализует преобразование в нескольких строках:
\begin{lstlisting}[language=Python,numbers=none]
P = 2**255 - 19
def montgomery_to_edwards(u_bytes, sign):
    u = int.from_bytes(u_bytes, "little") & ((1 << 255) - 1)
    y = (u - 1) * pow(u + 1, P - 2, P) % P      # деление с помощью обратного элемента по модулю p
    compressed = bytearray(y.to_bytes(32, "little"))
    compressed[31] |= (sign & 1) << 7          # знак координаты x в старшем бите
    return bytes(compressed)
\end{lstlisting}
Затем модуль проверяет перекрёстный сертификат с помощью декодера сертификатов из лабораторной работы~1. Применительно к дескрипторам лабораторной сети (ключ сокращён для удобства чтения; ключи меняются с каждой новой сетью):
\begin{lstlisting}[numbers=none]
$ PYTHONPATH=labs/ru/код python3 -m torlab.conversion labs/ru/лаборатория/net/nodes/000a/cached-descriptors test004r
test004r: ntor-onion-key DaxPE+Grnjex2Hmv+5qutl5dIX4La2jF6MIL7IK6+3E
  ключ Эдвардса (знак 0): c369d08f8083002a15fe22f65dca14bf...
  crosscert действителен; подтверждённая идентичность: kjPIgDB2OFsGJzuwsT+KMS/El9qrlE2/4P5WvIvmzC8
\end{lstlisting}
В приведённом примере подпись действительна, а подтверждённая идентичность совпадает с записанной в файле \file{fingerprint-ed25519} ретранслятора. Это же наблюдалось у остальных ретрансляторов лабораторной сети. При неверном знаке проверка завершается неудачей — такой случай предусмотрен отрицательным автоматизированным тестом модуля. Деление по модулю $p$ вычисляется с помощью малой теоремы Ферма: умножением на $(u+1)^{p-2}$.

% ---------------------------------------------------------------------------
\section{Изменяемость шифртекста потокового шифра}\label{sec:B-malleabilita}

AES в режиме счётчика (\cref{sec:cripto-cifrari}) создаёт поток байтов, который объединяется с сообщением операцией XOR. Изменив бит шифртекста, атакующий изменяет \emph{соответствующий бит открытого текста}, даже не зная ключа: достаточно знать или угадать исходное содержимое. Пример на Python:
\begin{lstlisting}[language=Python,numbers=none]
from cryptography.hazmat.primitives.ciphers import Cipher, algorithms, modes
import os
k, n = os.urandom(16), bytes(16)
cipher = lambda d: Cipher(algorithms.AES(k), modes.CTR(n)).encryptor().update(d)
c = bytearray(cipher(b"BEGIN www.example.org:443"))
delta = bytes(x ^ y for x, y in zip(b"www.example.org", b"www.example.net"))
for i, d in enumerate(delta):
    c[6 + i] ^= d                 # целевое изменение без знания ключа
print(cipher(bytes(c)))            # b'BEGIN www.example.net:443'
\end{lstlisting}
Исходное послойное шифрование ретрансляционных ячеек Tor (tor1) использует AES-CTR; целостность проверяется на конечном узле цепи по четырёхбайтовому дайджесту. Поэтому вредоносный ретранслятор способен предсказуемо изменить ячейку, а сообщник дальше по маршруту — распознать изменение. Это атака маркировки из \cref{sec:tor1-debolezze} и одна из причин создания CGO (\cref{sec:cgo}), где изменённая ячейка и последующие становятся невосстановимыми.

% ---------------------------------------------------------------------------
\section*{Упражнения}

\begin{esercizio}[1]\label{es:B-bit}
Сколько различных значений может принимать закрытый скаляр X25519 после фиксации битов? Сравните ответ с $2^{256}$.
\end{esercizio}

\begin{esercizio}[2]\label{es:B-inversa}
Выведите и реализуйте обратное преобразование от координаты $y$ открытого ключа Ed25519 к координате Монтгомери $u$. Сохраните исходный бит знака $x$, а затем с помощью функции \texttt{montgomery\_to\_edwards} из \file{conversion.py} проверьте, что обе операции восстанавливают исходный открытый ключ Ed25519 (при $y \neq 1$).
\end{esercizio}

\begin{esercizio}[3]\label{es:B-digest}
В tor1 ячейка содержит четырёхбайтовый дайджест и двухбайтовое нулевое поле \texttt{recognized}. Какова вероятность, что случайно изменённая ячейка будет принята последним узлом как допустимая? Почему атака маркировки всё же возможна?
\end{esercizio}
